\section*{Chapitre \ref{ch:ll:logging-capture-and-replay} --- Journalisation, capture et rejeu}

\begin{solution}{exo:ll:logging-capture-and-replay:1}
$16 + 4 \times 8 = 48$ octets. La ligne «~34200000000010 fill 40 of 100 at 999900, position -200~» compte 54 caractères, 55 avec son
retour à la ligne~: à peu près la même taille. Le \omterm{def:ll:logging-capture-and-replay:binlog}{journal binaire} économise le formatage (les conversions en décimal et l'analyse du format) sur
le chemin critique (hot path), et non des octets.
\end{solution}

\begin{solution}{exo:ll:logging-capture-and-replay:2}
Le décodeur lit les arguments avec les types enregistrés pour l'identifiant. Deux appels avec le même format et des
types différents (un \texttt{int} et un \texttt{double}) codent des octets différents~; s'ils partageaient un identifiant, l'un des
deux serait mal décodé.
\end{solution}

\begin{solution}{exo:ll:logging-capture-and-replay:3}
$12 + 189\,875 \times 56 = 10\,633\,012$ octets~; compressé, $2\,831\,089 / 189\,875 \approx 14.9$ octets par événement.
\end{solution}

\begin{solution}{exo:ll:logging-capture-and-replay:4}
$2 \times 10^6 \times 0.02 = 40\,000$ enregistrements arrivent pendant que l'écrivain est bloqué~; la capacité de l'anneau est une
puissance de deux, soit 65\,536 cases, 4~MiB avec des cases de 64 octets.
\end{solution}

\begin{solution}{exo:ll:logging-capture-and-replay:5}
Enregistrements d'étalonnage écrits régulièrement par l'écrivain (et non par le chemin critique)~: des couples formés d'une lecture du TSC et
d'une lecture d'une horloge traçable à l'UTC (disciplinée par PTP, chapitre~5), rapprochées dans le temps. Hors ligne, la valeur TSC de chaque
enregistrement est convertie par interpolation entre les couples d'étalonnage qui l'encadrent~; l'erreur est celle de l'horloge de référence augmentée de la dérive entre deux couples.
\end{solution}

\begin{solution}{exo:ll:logging-capture-and-replay:6}
Avec le \omterm{def:ll:logging-capture-and-replay:binlog}{journal binaire}, environ $5 \times \qty{5}{\nano\second} = \qty{25}{\nano\second}$ par événement, soit la moitié du travail propre à
l'événement~; avec \texttt{snprintf}, environ $5 \times \qty{100}{\nano\second} = \qty{500}{\nano\second}$, soit dix fois l'événement.
\end{solution}

\begin{solution}{exo:ll:logging-capture-and-replay:7}
Une instruction \texttt{log<\textquotedbl{}order \{\} refused \{\} mask \{\}\textquotedbl{}>(t, cl, code, mask)} dans le chemin de refus de la
barrière (les ordres acceptés peuvent avoir les leurs)~; rejouer \texttt{events.txt} à travers la barrière avec un journaliseur, vider le journal et
le décodage, les codes comptés par motif donnent les totaux de la première ligne de \texttt{expected.txt}.
\end{solution}

\begin{solution}{exo:ll:logging-capture-and-replay:8}
Un journal de décisions indique ce qui a été décidé, non ce qui l'a fondé, et il ne peut pas être rejoué sous un débogueur ni après une correction. L'enregistrement de la bourse
n'est pas ce que le processus a reçu~: il ne comporte ni trous, ni retransmissions, ni heures locales de réception, ni
comptes rendus d'ordres, ni changements de paramètres, ni déclenchements de minuteries, et il suit l'ordre de la bourse, non celui du processus. Un rejeu de cet
enregistrement répond à ce qu'aurait fait la stratégie avec des données parfaites, non à ce qu'elle a fait.
\end{solution}

\begin{solution}{pb:ll:logging-capture-and-replay:1}
\begin{enumerate}
\item $189\,875 / \qty{120}{\second} \approx 1\,582$ événements par seconde.
\item $1\,582 \times 6.5 \times 3\,600 \approx 37.0$ millions d'événements.
\item $37.0 \times 10^6 \times 56 \approx 2.07$~Go bruts, soit environ 0.55~Go compressés à 3.76 pour 1.
\item L'ouverture est l'un des moments les plus chargés de la journée, donc la moyenne calculée sur elle surestime un après-midi calme~; mais le stockage et
le rejeu doivent être dimensionnés pour les journées chargées, et les jours d'événements le sont davantage encore.
\item $0.55~\text{Go} \times 252 \times 5 \approx 0.7$~To.
\item Environ 139~To pour 200 instruments.
\item Le binaire exact et sa construction, sa configuration, les instantanés de paramètres, les données de référence (définitions des instruments),
les comptes rendus d'ordres qu'il a reçus, et les outils de rejeu, le tout versionné avec le journal.
\item Il contient les résultats, mais non leurs causes~: il ne peut être rejoué, arrêté à un événement, ni servir à tester une correction.
\item À environ 2\,800 fois le temps réel, $23\,400 / 2\,800 \approx \qty{8.4}{\second}$~: moins de dix secondes.
\item Une demi-heure environ, ou moins sur plusieurs cœurs, puisque les instruments se rejouent indépendamment.
\item La lecture et le décodage du journal, et le travail propre du moteur par événement (mémoire et processeur)~; jamais le rythme du marché,
puisque rien n'attend.
\item Toute décision dépendant du temps serait différente, et le rejeu devrait attendre que le temps réel s'écoule.
\item \textbf{Résultat chiffré.} Un journal coûte 56 octets par événement brut et environ 15 compressé (3.76 pour 1)~: à la moyenne de 1\,582 événements
par seconde de l'ouverture, une journée de 6.5 heures compte 37 millions d'événements, soit 2.1~Go bruts et 0.55~Go compressés par instrument, environ
0.7~To sur cinq ans~; et un rejeu va environ 2\,800 fois plus vite que le temps réel, une journée en moins de dix secondes.
\item Environ \qty{5}{\nano\second} avec le \omterm{def:ll:logging-capture-and-replay:binlog}{formatage différé}, environ \qty{100}{\nano\second} avec \texttt{snprintf}.
\item Les mêmes actions, bit à bit, jusqu'à l'ordre en question inclus, avec l'état qui l'a produit~; vérifié en
comparant le journal de décisions rejoué (ou son empreinte) au journal de production.
\item Elle est perdue, à moins que les lectures d'horloge ne soient elles-mêmes journalisées comme entrées et rejouées.
\item Sur le fil du moteur, là où les entrées y pénètrent, dans l'ordre où il les traite~: plus tôt, l'ordre pourrait
différer~; plus tard, quelque chose pourrait manquer.
\item Que la comparaison de deux journaux trouve la première divergence et que rien avant elle ne diffère~; pour un cas réel, la même
comparaison entre un journal de production et son rejeu, puis l'événement du journal à ce point, examiné dans le rejeu.
\item Une microseconde ou mieux, à moins de 100 microsecondes de l'UTC pour le trading à haute fréquence (HFT)~; la résolution du TSC est bien plus fine,
et avec un étalonnage contre une horloge traçable à l'UTC, le journal satisfait aux deux.
\item Un journal complet de ses entrées, un programme déterministe et la construction exacte qui l'a exécuté.
\end{enumerate}
\end{solution}

\begin{solution}{iq:ll:logging-capture-and-replay:1}
Elle met en forme du texte, analyse un format, peut prendre un verrou, allouer ou écrire dans un fichier, et sa
\omterm{def:ll:where-latency-comes-from:tail}{latence de queue} est longue~: cent nanosecondes ou plus par ligne. Écrire un identifiant d'instruction, un horodatage et les arguments bruts dans un anneau, laisser un autre fil les écrire, et
formater hors ligne.
\iqlookfor{enregistrements binaires et \omterm{def:ll:logging-capture-and-replay:binlog}{formatage différé}.}
\end{solution}

\begin{solution}{iq:ll:logging-capture-and-replay:2}
Des identifiants fixés à la compilation, un enregistrement de taille fixe encodé par quelques copies dans un anneau mono-producteur préalloué par
fil, un fil d'écriture sur un autre cœur qui vide l'anneau par lots. Quand l'anneau est plein, abandonner l'enregistrement et le compter~: le chemin
critique n'attend jamais, et le compte dit ce qui a été perdu.
\iqlookfor{aucun travail au-delà de la copie, et une politique d'abandon.}
\end{solution}

\begin{solution}{iq:ll:logging-capture-and-replay:3}
Chaque entrée telle que reçue (données de marché après le \omterm{def:ll:the-feed-handler:handler}{feed handler}, comptes rendus d'ordres, minuteries dépendant du temps, changements
de paramètres, messages de contrôle) avec son heure de réception, plus la construction et la configuration~; enregistrées sur le fil du moteur là
où les entrées y pénètrent, dans leur ordre.
\iqlookfor{entrées telles que reçues, heure comprise, à la frontière du moteur.}
\end{solution}

\begin{solution}{iq:ll:logging-capture-and-replay:4}
Journaliser les sorties des deux avec assez de détail (actions, ou empreintes de l'état à chaque événement), les comparer enregistrement par
enregistrement, et aller à la première différence~; regarder ensuite ce qui différait à cet événement~: temps, ordre d'itération, fils, entrées.
\iqlookfor{la première différence, puis la cause.}
\end{solution}

\begin{solution}{iq:ll:logging-capture-and-replay:5}
Le journal de la journée, la construction exacte, sa configuration et ses paramètres, et les outils de rejeu~; rejouer jusqu'à l'ordre, vérifier le
rejeu contre le journal de production, et lire l'état qui a produit la décision. Avec les journaux et les outils conservés, cela prend
quelques minutes.
\iqlookfor{journal plus construction, et rejeu vérifié.}
\end{solution}

\begin{solution}{iq:ll:logging-capture-and-replay:6}
Chaque processus journalise ses propres entrées dans un format binaire commun, avec des horodatages traçables à l'UTC, et journalise ses décisions en
binaire, et envoie les deux vers le stockage, compressés~; les constructions et les configurations sont versionnées avec eux~; des pilotes de rejeu
par composant, tournent en intégration continue sur les journées récentes et comparent les résultats aux journaux de production~; la rétention suit la réglementation.
\iqlookfor{journaux par processus, format commun, versionnement et rejeu en CI.}
\end{solution}
